\documentclass[11pt]{article}
\usepackage{amsmath,amssymb,amsthm}
\usepackage[margin=1.1in]{geometry}
\usepackage{booktabs}
\usepackage[hidelinks]{hyperref}

\newtheorem{theorem}{Theorem}
\newtheorem{proposition}[theorem]{Proposition}
\newtheorem{lemma}[theorem]{Lemma}
\newtheorem{corollary}[theorem]{Corollary}
\theoremstyle{definition}
\newtheorem{question}[theorem]{Question}
\newtheorem{remark}[theorem]{Remark}

\newcommand{\SE}{\mathrm{SE}}
\newcommand{\N}{\mathbb{N}}
\DeclareMathOperator{\modstar}{mod^{*}}

\title{A negative answer to a question of Kierstead, Kostochka and Xiang\\
on strongly equitable list coloring of forests}
\author{}
\date{September 2026 --- draft}

\begin{document}
\maketitle

\begin{abstract}
A forest $T$ on $n$ vertices is equitably $k$-colorable ($k\ge 3$) if and only if every
vertex lies in an independent set of size at least $\lfloor n/k\rfloor$ (Chen--Lih;
Miyata--Tokunaga--Kaneko). Kierstead, Kostochka and Xiang asked whether the same
condition characterizes strongly equitable $k$-choosability of forests. We show that it
does not: for every $k\ge 2$ there is a forest, and even a tree, on $2k^3+k$ vertices
that satisfies the independence condition (indeed is equitably $k$-colorable) but is not
equitably $k$-choosable even in the weaker sense of Kostochka, Pelsmajer and West.
The construction sharpens an example of Kaul, Mudrock and Wagstrom, and it generalizes
to a family of necessary conditions for strongly equitable choosability of forests.
\end{abstract}

\section{Introduction}

All graphs are finite and simple. For a graph $G$ and a vertex $v$, let $\alpha_v(G)$
denote the maximum size of an independent set of $G$ containing $v$, and let $N(v)$,
$N[v]=N(v)\cup\{v\}$ and $d(v)=|N(v)|$ be the neighborhood, closed neighborhood and
degree of $v$. We write $[k]=\{1,\dots,k\}$.

A proper vertex coloring is \emph{equitable} if any two color classes differ in size by
at most one. In an equitable $k$-coloring of an $n$-vertex graph every class has at least
$\lfloor n/k\rfloor$ vertices, so the class of $v$ witnesses
$\alpha_v(G)\ge\lfloor n/k\rfloor$. For forests this obvious necessary condition is also
sufficient.

\begin{theorem}[Chen--Lih \cite{CL}; Miyata--Tokunaga--Kaneko \cite{MTK}; see Chang \cite{Chang}]
\label{thm:CL}
Let $k\ge 3$. A forest $T$ on $n$ vertices is equitably $k$-colorable if and only if
$\alpha_v(T)\ge\lfloor n/k\rfloor$ for every vertex $v$.
\end{theorem}

We call the condition in Theorem~\ref{thm:CL} the \emph{independence condition}.

\subsection*{List versions}
A \emph{$k$-list assignment} for $G$ is a map $L$ from $V(G)$ to sets of colors with
$|L(v)|=k$ for all $v$, and an \emph{$L$-coloring} is a proper coloring $f$ with
$f(v)\in L(v)$ for all $v$. Kostochka, Pelsmajer and West \cite{KPW} call an $L$-coloring
\emph{equitable} if every color class has at most $\lceil |G|/k\rceil$ vertices, and $G$
\emph{equitably $k$-choosable} if it has an equitable $L$-coloring for every $k$-list
assignment $L$. Note that an equitable $L$-coloring may use more than $k$ colors and may
have very small classes; only an upper bound on class sizes is imposed.

Kierstead, Kostochka and Xiang \cite{KKXplanar,KKXsparse} introduced a stronger notion.
Write $|G|=kq+r$ with $q\in\N$ and $1\le r\le k$, and set $|G|\modstar k:=r$ (so
$|G|\modstar k=k$ when $k$ divides $|G|$). A color class is \emph{full} if it has exactly
$\lceil |G|/k\rceil$ vertices. An $L$-coloring is \emph{strongly equitable} ($\SE$) if no
class has more than $\lceil |G|/k\rceil$ vertices and at most $|G|\modstar k$ classes are
full; $G$ is \emph{$\SE$ $k$-choosable} if it has an $\SE$ $L$-coloring for every $k$-list
assignment $L$. Two features of this definition matter here.
\begin{itemize}
\item Every $\SE$ $L$-coloring is an equitable $L$-coloring in the sense of \cite{KPW};
  hence $\SE$ $k$-choosable graphs are equitably $k$-choosable.
\item If all lists equal $[k]$, then the $\SE$ $L$-colorings are exactly the equitable
  $k$-colorings: with $r=|G|\modstar k$ there are $r$ classes of size $q+1$ and $k-r$ of
  size $q$. Hence $\SE$ $k$-choosable graphs are equitably $k$-colorable, and in particular
  $\SE$ $k$-choosable forests satisfy the independence condition.
\end{itemize}
In their survey \cite[Question~13]{KKXsurvey}, Kierstead, Kostochka and Xiang asked
whether Theorem~\ref{thm:CL} extends to this setting.

\begin{question}[\cite{KKXsurvey}]\label{q13}
Is it true that for each $k\ge 3$, an $n$-vertex forest $T$ is $\SE$ $k$-choosable if and
only if $\alpha_v(T)\ge\lfloor n/k\rfloor$ for every vertex $v$?
\end{question}

By the second feature above, the ``only if'' direction holds. We show that the ``if''
direction fails for every $k$, even for trees, and even for the weaker notion of
\cite{KPW}.

\begin{theorem}\label{thm:main}
Let $k\ge 2$, $q=2k^2+1$ and $n=kq=2k^3+k$. There is a forest $F_k$ and a tree $T_k$, each
on $n$ vertices, such that
\begin{enumerate}
\item[\textup{(i)}] $F_k$ and $T_k$ have equitable $k$-colorings; in particular
  $\alpha_v\ge q=\lfloor n/k\rfloor$ for every vertex $v$;
\item[\textup{(ii)}] $F_k$ and $T_k$ are not equitably $k$-choosable in the sense of
  \cite{KPW}, and hence not $\SE$ $k$-choosable.
\end{enumerate}
Consequently the answer to Question~\ref{q13} is negative for every $k\ge 3$.
\end{theorem}

The forest $F_k$ is a disjoint union of two stars. Kaul, Mudrock and Wagstrom
\cite[Proposition~22]{KMW} proved, for a different purpose, that
$K_{1,(k-1)(k^3-k+2)}+K_{1,k^3}$ is not equitably $k$-choosable. That forest also satisfies
the independence condition, so it already answers Question~\ref{q13} in the negative;
the connection does not seem to have been noticed. Our $F_k$ is the smallest member of
the same family of examples (Remark~\ref{rem:KMW}). In Section~\ref{sec:general}
the argument is generalized to a family of necessary conditions for $\SE$
$k$-choosability of forests, and Section~\ref{sec:open} collects open questions.

\section{The construction}\label{sec:construction}

Throughout this section $k\ge 2$, $q=2k^2+1$ and $n=kq$. Since $k$ divides $n$, we have
$\lfloor n/k\rfloor=\lceil n/k\rceil=q$, and an $L$-coloring is equitable in the sense of
\cite{KPW} exactly when every color is used at most $q$ times.

\subsection*{The forest}
Let $F_k$ be the disjoint union of two stars: one with center $u_0$ and leaf set $A$,
where $|A|=(k-1)q-1$, and one with center $w_0$ and leaf set $B$, where $|B|=2k^2$.
Then
\[
|F_k| = 2+\big((k-1)q-1\big)+2k^2 = (k-1)q+q = kq = n .
\]
Index the leaves of $w_0$ by triples,
$B=\{b_{c,d,t} : c\in[k],\ d\in\{k+1,\dots,2k\},\ t\in\{1,2\}\}$.

\subsection*{Equitable colorability}
The set $\{u_0\}\cup B$ is independent and has $1+2k^2=q$ vertices; the set
$\{w_0\}\cup A$ is independent and has $(k-1)q$ vertices. Using $\{u_0\}\cup B$ as one
class and splitting $\{w_0\}\cup A$ into $k-1$ classes of size $q$ gives an equitable
$k$-coloring of $F_k$. This proves (i) of Theorem~\ref{thm:main} for $F_k$.

\subsection*{The list assignment}
Define a $k$-list assignment $L$ by
\[
L(u_0)=L(a)=[k]\ \ (a\in A),\qquad
L(w_0)=\{k+1,\dots,2k\},\qquad
L(b_{c,d,t})=\big([k]\setminus\{c\}\big)\cup\{d\}.
\]

\subsection*{No equitable $L$-coloring}
Suppose $f$ is an $L$-coloring of $F_k$ in which every color is used at most $q$ times.
Let $c=f(u_0)\in[k]$ and $d=f(w_0)\in\{k+1,\dots,2k\}$. Consider the set
\[
S = A\cup\{b_{c,d,1},\,b_{c,d,2}\},\qquad |S|=(k-1)q+1 .
\]
Every $a\in A$ is adjacent to $u_0$ and has list $[k]$, so $f(a)\in[k]\setminus\{c\}$.
Each of $b_{c,d,1},b_{c,d,2}$ is adjacent to $w_0$, so it does not receive $d$, and its
remaining colors are $[k]\setminus\{c\}$. Thus $f$ maps the $(k-1)q+1$ vertices of $S$
into the $k-1$ colors of $[k]\setminus\{c\}$, and by the pigeonhole principle some color
is used more than $q$ times, a contradiction. Hence $F_k$ is not equitably $k$-choosable,
and, since every $\SE$ $L$-coloring is an equitable $L$-coloring, not $\SE$
$k$-choosable. This proves (ii) for $F_k$.

\subsection*{The tree}
Fix $a_0\in A$ and $b_0\in B$ and let $T_k=F_k+a_0b_0$. Adding an edge cannot create an
$L$-coloring, so $T_k$ is not equitably $k$-choosable. For (i), $\{u_0\}\cup B$ and
$\{w_0\}\cup A$ are still independent in $T_k$ (neither contains both $a_0$ and $b_0$), and
the equitable $k$-coloring above places $a_0$ and $b_0$ in different classes, so it is a
proper equitable $k$-coloring of $T_k$. This completes the proof of
Theorem~\ref{thm:main}. \qed

\begin{remark}\label{rem:KMW}
In the construction, the leaves of $u_0$ fill the $k-1$ colors $[k]\setminus\{c\}$ up to a
\emph{slack} of exactly one vertex, and for each of the $k^2$ possible pairs
$(f(u_0),f(w_0))$ two leaves of $w_0$ compete for that single slot. The example of
\cite[Proposition~22]{KMW} is the same mechanism with slack $k-1$ and $k$ competing
leaves per pair, which costs $n=k(k^3-k+3)$ vertices instead of $2k^3+k$
($81$ instead of $57$ for $k=3$). Slack zero is impossible: if $|A|=(k-1)q$ then
$|B|=q-2$ and $\alpha_{u_0}=q-1$, which violates the independence condition. This is
precisely the obstruction that makes Theorem~\ref{thm:CL} true, and the list version
escapes it by one vertex.
\end{remark}

\section{A general obstruction}\label{sec:general}

The proof of Theorem~\ref{thm:main} only used a count of vertices forced into $k-1$
colors. We record the general form of this count.

Let $T$ be a forest on $n$ vertices and write $n=kq+r'$ with $0\le r'\le k-1$, so that
$q=\lfloor n/k\rfloor$. In an $\SE$ $L$-coloring, any $k-1$ colors together carry at most
$n-q$ vertices: if $r'=0$ every class has at most $q$ vertices, and if $r'\ge 1$ every
class has at most $q+1$ vertices with at most $r'$ classes of size $q+1$, so the bound is
$(k-1)q+r'=n-q$ in both cases. In an equitable $L$-coloring in the sense of \cite{KPW}
the corresponding bound is $(k-1)\lceil n/k\rceil$.

\begin{proposition}\label{prop:obstruction}
Let $T$ be a forest on $n$ vertices, $q=\lfloor n/k\rfloor$, let $u\in V(T)$, let
$W\subseteq V(T)\setminus\{u\}$, and let $Y_w\subseteq N(w)\setminus N[u]$ $(w\in W)$ be
pairwise disjoint. If
\begin{equation}\label{eq:obstruction}
d(u)-|W\cap N(u)|+\sum_{w\in W}\Big\lfloor \frac{|Y_w|}{k^2}\Big\rfloor \;>\; n-q ,
\end{equation}
then $T$ is not $\SE$ $k$-choosable. If \eqref{eq:obstruction} holds with $n-q$ replaced
by $(k-1)\lceil n/k\rceil$, then $T$ is not equitably $k$-choosable in the sense of
\cite{KPW}.
\end{proposition}

\begin{proof}
Let $D=\{k+1,\dots,2k\}$. Define $L(u)=[k]$, $L(x)=[k]$ for $x\in N(u)\setminus W$, and
$L(w)=D$ for $w\in W$. For each $w\in W$ split $Y_w$ into $k^2$ parts
$Y_{w,c,d}$ ($c\in[k]$, $d\in D$), each of size at least $\lfloor |Y_w|/k^2\rfloor$, and
give every $y\in Y_{w,c,d}$ the list $([k]\setminus\{c\})\cup\{d\}$. All other lists are
arbitrary $k$-sets. Let $f$ be an $\SE$ $L$-coloring and put $c=f(u)$. Every
$x\in N(u)\setminus W$ is colored from $[k]\setminus\{c\}$. For $w\in W$ let $d=f(w)\in D$;
every $y\in Y_{w,c,d}$ is adjacent to $w$, cannot receive $d$, and is therefore colored
from $[k]\setminus\{c\}$. The number of these vertices is at least the left-hand side of
\eqref{eq:obstruction}, which exceeds the capacity $n-q$ of $k-1$ colors. The second
statement is identical with the other capacity.
\end{proof}

Theorem~\ref{thm:main} is the case $W=\{w_0\}$, $Y_{w_0}=B$, $|B|=2k^2$,
$d(u_0)=(k-1)q-1$, where the left-hand side equals $(k-1)q+1=n-q+1$.
With $W=\emptyset$, condition \eqref{eq:obstruction} reads $d(u)>n-q$; since
$\alpha_u\le n-d(u)$, this is already excluded by the independence condition. Thus
Proposition~\ref{prop:obstruction} adds new information only when some vertex $w$ has at
least $k^2$ neighbors outside $N[u]$.

\section{Concluding remarks and open questions}\label{sec:open}

\subsection*{Verification}
The forest $F_k$, the list assignment $L$, the independence condition and the
non-existence of an $\SE$ $L$-coloring have been formalized for all $k\ge 2$ and
machine-checked in Lean~4 with Mathlib (file \texttt{Q13lean/Counterexample.lean} in the
accompanying repository). Independently, an exact flow-based program confirms
Theorem~\ref{thm:main} for $k=3,4$ and the example of \cite{KMW} for $k=3,4$.

\subsection*{Small forests}
An exhaustive computer search found no forest satisfying the independence condition and
admitting a defeating $k$-list assignment in the following ranges, where the palette is
the set of colors from which all lists are drawn.
\begin{center}
\begin{tabular}{@{}ccc@{}}
\toprule
$k$ & number of vertices & palette size\\
\midrule
$3$ & $\le 12$ & $4$\\
$3$ & $\le 10$ & $5$\\
$4$ & $\le 11$ & $5$\\
$5$ & $\le 10$ & $6$\\
\bottomrule
\end{tabular}
\end{center}
Because a defeating assignment may need a larger palette ($L$ above uses $2k$ colors),
these searches do not yield a lower bound on the size of a smallest counterexample.

\begin{question}
For $k=3$, what is the smallest number of vertices of a forest that satisfies the
independence condition but is not $\SE$ $3$-choosable? It is at most $57$.
\end{question}

\subsection*{A corrected characterization?}
The independence condition together with the family of conditions
``\eqref{eq:obstruction} fails for all choices of $u$, $W$, $(Y_w)$'' is necessary for
$\SE$ $k$-choosability of a forest. We do not know whether it is sufficient. A
first test case is the class of disjoint unions of two stars, where even equitable
$k$-choosability in the sense of \cite{KPW} has not been characterized \cite[Question~8]{KMW}.

\begin{question}
Is a forest $T$ on $n$ vertices $\SE$ $k$-choosable whenever $\alpha_v(T)\ge\lfloor n/k\rfloor$
for all $v$ and \eqref{eq:obstruction} fails for all $u$, $W$ and $(Y_w)$?
In particular, is $T$ $\SE$ $k$-choosable whenever the independence condition holds and
$\Delta(T)< k^2$?
\end{question}

For comparison, Kostochka, Pelsmajer and West \cite{KPW} proved that every forest is
equitably $k$-choosable (in the sense of \cite{KPW}) when $k\ge 1+\Delta/2$, and that
this bound is tight (as quoted in \cite{KMW}). No such degree bound can hold for $\SE$
$k$-choosability without the independence condition, since $\SE$ $k$-choosable graphs
are equitably $k$-colorable; whether the independence condition suffices for
$\SE$ $k$-choosability of forests with $\Delta\le 2k-2$ is a natural special case of
the question above. By Proposition~\ref{prop:obstruction}, the new list obstructions
appear only at degree $k^2$ and beyond.

\begin{thebibliography}{9}
\bibitem{Chang} G.~J.~Chang, A note on equitable colorings of forests,
  \emph{European J.~Combin.} 30 (2009) 809--812.
\bibitem{CL} B.-L.~Chen, K.-W.~Lih, Equitable coloring of trees,
  \emph{J.~Combin.\ Theory Ser.~B} 61 (1994) 83--87.
\bibitem{KMW} H.~Kaul, J.~A.~Mudrock, T.~Wagstrom, On the equitable choosability of the
  disjoint union of stars, arXiv:2008.06333 (2020).
\bibitem{KKXplanar} H.~A.~Kierstead, A.~V.~Kostochka, Z.~Xiang, Equitable list coloring
  of planar graphs with given maximum degree, \emph{J.~Graph Theory} 108 (2024) 832--838.
\bibitem{KKXsparse} H.~A.~Kierstead, A.~V.~Kostochka, Z.~Xiang, Equitable list coloring
  of sparse graphs, \emph{Discrete Math.} 349 (2026) 114981; arXiv:2411.08372.
\bibitem{KKXsurvey} H.~A.~Kierstead, A.~V.~Kostochka, Z.~Xiang, Results and problems on
  equitable coloring of graphs, in: \emph{Sum(m)it 280: Surveys in Extremal Combinatorics
  and Combinatorial Geometry}, Bolyai Soc.\ Math.\ Stud.\ 32 (2026); arXiv:2504.14711.
\bibitem{KPW} A.~V.~Kostochka, M.~J.~Pelsmajer, D.~B.~West, A list analogue of equitable
  coloring, \emph{J.~Graph Theory} 44 (2003) 166--177.
\bibitem{MTK} D.~Miyata, S.~Tokunaga, A.~Kaneko, Equitable coloring of trees, unpublished
  manuscript (cited in \cite{Chang,KKXsurvey}).
\end{thebibliography}

\end{document}
